% =====================================================================
%  Solver comparison for the Laplace-Beltrami problem on a torus
%
%  Drop-in section for the thesis. Requires only:
%      \usepackage{booktabs}
%      \usepackage{graphicx}
%  with scaling.pdf placed next to the document (or on the graphics path).
%
%  Every number here is taken from bench/results-d3.csv, produced by
%  bench/run_sweep.sh and analysed by bench/analyse.py. Do not edit the
%  numbers by hand; regenerate instead.
% =====================================================================

\section{Solver comparison: multigrid against its alternatives}
\label{sec:solver-comparison}

The geometric multigrid method of Section~\ref{sec:multigrid} is not the only
way to solve the discrete system. This section places it against the standard
alternatives for a symmetric positive definite operator, and shows that its
advantage is not a favourable constant but a different asymptotic behaviour
under mesh refinement.

\subsection{What was compared}

Five methods were applied to the \emph{same} linear system: one torus, cubic
Lagrange elements, one curved mapping, one right-hand side. Every method
therefore reproduces the same discretisation error, and the comparison isolates
the solver. Each run was performed in a separate process, and the multigrid
level hierarchy is assembled only for the methods that use it, so no method is
charged for machinery it does not touch. Two of the five are the same V-cycle:
in Table~\ref{tab:methods}, the multigrid rows differ only in whether a Krylov
subspace steers the iteration.

\begin{table}[htbp]
  \centering
  \caption{The five methods compared. All iterative methods stop at a relative
    residual of $10^{-10}$.}
  \label{tab:methods}
  \begin{tabular}{@{}ll@{}}
    \toprule
    Method & Description \\
    \midrule
    CG $+$ Jacobi      & conjugate gradients, diagonal preconditioner \\
    CG $+$ SSOR        & conjugate gradients, symmetric SOR preconditioner \\
    MG standalone      & the V-cycle iterated on the defect, no Krylov acceleration \\
    CG $+$ MG          & the same V-cycle used as a preconditioner for CG \\
    Direct             & sparse LU factorisation (UMFPACK) \\
    \bottomrule
  \end{tabular}
\end{table}

Two memory quantities are reported throughout, because they answer different
questions. \emph{Objects} counts the matrices, vectors and degree-of-freedom
structures the solver holds directly. \emph{Process} is the peak resident
memory above a measured baseline. The difference between the two is memory
allocated inside a library, which is precisely where a direct solver's fill-in
lives and where the first quantity sees nothing.

\subsection{Results at the largest problem size}

Table~\ref{tab:largest} gives one row per method at $N = 589{,}824$ unknowns
(65{,}536 cells, seven uniform refinement cycles) --- the largest problem the
study could run, for the reason given below. \emph{Linear algebra} is the
preconditioner or factorisation together with its application, excluding
assembly: that work is the same $\mathcal{O}(N)$ cost for every method and
would otherwise dilute the comparison.

\begin{table}[htbp]
  \centering
  \caption{All five methods at $N = 589{,}824$ unknowns. The linear-algebra
    column excludes assembly, which is common to all five. Memory is given
    both as solver objects and as process memory above baseline: the two
    answer different questions, and the gap between them is memory allocated
    inside a library.}
  \label{tab:largest}
  \small
  \begin{tabular}{@{}lrrrrr@{}}
    \toprule
    Method & Iterations & Lin.\ alg.\ (s) & Total (s)
           & Objects (MB) & Process (MB) \\
    \midrule
    CG $+$ Jacobi      & 1000\textsuperscript{$\dagger$} & 179.1 & 309.2 & 203.5 & \phantom{0}430.5 \\
    CG $+$ SSOR        & \phantom{0}957 & 692.6 & 835.6 & 203.5 & \phantom{0}451.7 \\
    MG standalone      & \phantom{00}31 & 114.3 & 335.6 & 443.5 & \phantom{0}833.0 \\
    CG $+$ MG          & \phantom{00}15 & \phantom{0}62.1 & 329.7 & 443.5 & \phantom{0}793.1 \\
    Direct (UMFPACK)   & \phantom{000}1 & 652.0 & 800.4 & 203.5 & 3347.3 \\
    \bottomrule
    \multicolumn{6}{@{}l@{}}{\footnotesize\textsuperscript{$\dagger$}Exhausted its budget of
      1000 iterations without reaching tolerance.}
  \end{tabular}
\end{table}

Three observations follow.

\paragraph{Jacobi does not solve this problem at all.} It exhausts its budget
of 1000 iterations with the relative residual still at $9.7 \times 10^{-8}$,
three orders of magnitude short of what every other method attains. Its
$H^1$ error, $2.570 \times 10^{-7}$, is consequently \emph{worse} than the
other four, which agree with one another at $2.549 \times 10^{-7}$. This is
the only point in the sweep at which the methods cease to produce the same
answer, and it is why the benchmark records non-convergence as a result rather
than treating it as a failure. Its row also appears inexpensive at $309$~s
only because it stopped early without solving the problem.

\paragraph{The direct solver requires 3.3~GB to factor a 204~MB matrix.} That
is a factor of $16.4$. Its \emph{object} memory is unremarkable --- the same
$203.5$~MB held by Jacobi and SSOR, the active matrix and nothing more --- but
the \emph{process} memory reveals what the first measure cannot: fill-in and
its internal index arrays, allocated inside UMFPACK. The growth is not
academic one level further on. Over the final refinement its memory rose by a
factor of $4.0$ for a factor of $4$ in unknowns, so the next refinement
projects beyond the $8$~GB available on the machine used here.
\emph{$589{,}824$ unknowns is therefore a ceiling imposed by memory, not a
stopping point chosen for convenience.}

\paragraph{Multigrid is the fastest method as well as the most economical in
the regime that matters.} It is the only one of the five whose advantage is
structural rather than a constant factor, which is the subject of the next
subsection.

\subsection{The scaling is the argument}

A single problem size proves little, because a direct factorisation has a small
constant and is competitive while the problem is small. What distinguishes the
methods is how each responds to refinement. Table~\ref{tab:exponents} reports
the exponent $b$ in $\log y = a + b\log N$, fitted over the largest half of the
converged points, and Figure~\ref{fig:scaling} shows the same data.

\begin{table}[htbp]
  \centering
  \caption{Observed scaling exponents, fitted on the largest half of the
    converged points. A method with $b = 1$ for time and memory is optimal;
    larger $b$ means the method degrades under refinement.}
  \label{tab:exponents}
  \small
  \begin{tabular}{@{}lrrrr@{}}
    \toprule
    Method & Iterations & Linear algebra & Total time & Process memory \\
    \midrule
    CG $+$ Jacobi    & $N^{0.50}$ & $N^{1.26}$ & $N^{1.08}$ & $N^{0.51}$ \\
    CG $+$ SSOR      & $N^{0.47}$ & $N^{1.40}$ & $N^{1.28}$ & $N^{0.66}$ \\
    MG standalone    & $N^{0.00}$ & $N^{0.89}$ & $N^{0.94}$ & $N^{0.74}$ \\
    CG $+$ MG        & $N^{0.02}$ & $N^{0.90}$ & $N^{0.99}$ & $N^{0.74}$ \\
    Direct (UMFPACK) & $N^{0.00}$ & $N^{1.40}$ & $N^{1.26}$ & $N^{0.93}$ \\
    \bottomrule
  \end{tabular}
\end{table}

\begin{figure}[htbp]
  \centering
  \includegraphics[width=\textwidth]{scaling.pdf}
  \caption{Scaling of the five methods from $144$ to $589{,}824$ unknowns,
    log--log in both axes, so a straight line is a power law and its slope is
    the exponent. Top left: iterations to reach tolerance. Top right: time in
    the linear algebra. Bottom left: memory held as solver objects. Bottom
    right: process memory above baseline. Hollow markers denote runs that
    exhausted their iteration budget and are therefore lower bounds; they are
    excluded from the fitted exponents.}
  \label{fig:scaling}
\end{figure}

\paragraph{Only multigrid has a mesh-independent iteration count.} Jacobi
requires 15 iterations at 144 unknowns and 673 at 147{,}456 --- the count
doubles with every refinement, $N^{0.50}$, which is the classical
$\mathcal{O}(h^{-1})$ growth of an unpreconditioned Krylov method; by
589{,}824 unknowns it has ceased to converge within the budget at all.
Multigrid's count moves from 12 to 15 across the entire range, $N^{0.02}$. It
is not that multigrid's iterations are cheaper; it is that there are always
approximately the same number of them. This is the property that makes the
method valuable, and it is visible in the top-left panel of
Figure~\ref{fig:scaling} as a horizontal line where every other method climbs.

\paragraph{The consequence appears in the time exponent.} Because Jacobi's
iteration count grows while each iteration costs $\mathcal{O}(N)$, its linear
algebra grows like $N^{1.26}$; SSOR is worse at $N^{1.40}$. Multigrid grows
like $N^{0.90}$. At the largest size this is the difference between $62$~s and
$693$~s, and the gap widens without bound, because these are exponents and not
constants.

\paragraph{Memory is where the direct solver is defeated, and the matrix does
not show it.} The direct method's object memory tracks the others exactly; its
process memory grows like $N^{0.93}$ against multigrid's $N^{0.74}$, from a
much larger constant. Multigrid does hold more than the simple
preconditioners --- $443.5$~MB against $203.5$~MB, because it stores an
operator for every level --- but this is a bounded factor, roughly $4/3$ of the
fine-grid matrix in two dimensions, and the two multigrid curves sit on lines
parallel to the others rather than steeper. A constant factor is a price; a
different slope is a ceiling.

\subsection{Multigrid as a solver and as a preconditioner}
\label{sec:mg-solver-vs-preconditioner}

The two multigrid variants are implemented on the same hierarchy and the same
V-cycle, so they differ in exactly one respect: whether the iteration is plain
defect correction or is accelerated by conjugate gradients. The result is
worth stating carefully, because it is less dramatic than it is sometimes
presented.

Both are mesh-independent, at $N^{0.00}$ and $N^{0.02}$ respectively. Standalone
multigrid therefore already delivers the property that motivates the method;
it does not require Krylov acceleration in order to be an optimal-order solver.
What conjugate gradients contributes is a factor of approximately two: it
requires 15 applications of the V-cycle where defect correction requires 31,
and that ratio holds at every problem size in the sweep.

This is the expected outcome. Defect correction is a fixed-point iteration
converging at the rate $\rho$ of the V-cycle, whereas conjugate gradients
minimises over a growing Krylov subspace and converges at a rate governed by
$\sqrt{\kappa}$ instead. Conjugate gradients additionally guarantees monotone
decrease of the error in the energy norm, which a bare V-cycle iteration does
not, and it degrades gracefully if the smoother or the coarse-grid solve is
imperfect. The honest summary is that CG $+$ MG is about twice as fast, not
asymptotically better.

\subsection{Limitations of these measurements}
\label{sec:limitations}

The timings reported above were taken inside an \texttt{arm64} container
running under emulation on an \texttt{x86\_64} host, because that was the only
available build of the finite element library. The iteration counts, the
residuals, the discretisation errors and the memory in bytes are exact and
independent of the architecture. \emph{The wall-clock times are not.} The
fitted exponents are ratios and largely survive the distortion, but the
absolute times in Table~\ref{tab:largest} should not be quoted as
hardware-representative; a native build would be required for figures intended
for direct comparison with published timings.

Two further caveats apply. First, assembly is excluded from the
linear-algebra column but not from the totals, and at the largest size it
dominates --- 162~s of multigrid's 330~s total --- so differences in total time
understate the differences between the solvers themselves. Second, all results
are for a single operator with a smooth variable coefficient,
$\kappa(\xi,\eta) = 1.1 + \sin^2\xi\cos^2\eta$, bounded away from zero.
Multigrid's advantage is sensitive to that assumption; a coefficient with large
jumps would require a different smoother and a different coarse-grid
correction, and the exponents above would not be expected to carry over
unchanged.
